\documentclass[10pt]{article}
\usepackage[margin=0.72in]{geometry}
\usepackage{booktabs}
\usepackage{graphicx}
\usepackage{float}
\usepackage[hidelinks]{hyperref}
\hypersetup{pdftitle={Spine Formal-v7 Primary Dynamic-Graph Results}}
\title{\textbf{Spine Formal-v7 Primary Dynamic-Graph Results}\\
\large Correctness-Gated Spine vs. GraSU+ReGraph K4-shared}
\author{Chuxiao Han}
\date{Live snapshot, July 2026}
\IfFileExists{data/formal_v7_primary/pair_table.tex}{
  \newcommand{\vdatadir}{data/formal_v7_primary}
  \newcommand{\rqdatadir}{data/rq3}
  \newcommand{\vfigdir}{../figures}
}{
  \newcommand{\vdatadir}{docs/paper/data/formal_v7_primary}
  \newcommand{\rqdatadir}{docs/paper/data/rq3}
  \newcommand{\vfigdir}{docs/figures}
}
\begin{document}
\maketitle
\begin{abstract}
This live report contains 25 of
46 expected formal-v7 executions and
9 complete Spine/K4-shared pairs. Missing bars are unfinished,
policy-stopped, or behavior-transition-invalidated executions awaiting an
identical-case successor, never zero-valued measurements. Every admitted row
passes architecture-precision and independent mathematical oracles, full
final-state comparison, and request, response, and DRAM conservation.
\end{abstract}

\section{Measurement contract}
\begin{table}[H]
\centering
\small
\input{\vdatadir/measurement_table.tex}
\caption{Measurement windows observed in admitted formal-v7 rows.}
\end{table}

\paragraph{Dataset scope.}
The frozen corpus contains 11 real datasets;
9 enter the full-graph comparison. The
remaining datasets, BC (24,575,382 vertices), UK (18,520,486 vertices), exceed the frozen Spine capacity of
16,777,216 vertices. They are capacity exclusions,
not failed or selectively removed performance rows, and no slices replace them
in the full-graph aggregate.

\paragraph{Synthetic endpoint.}
R19-32 (524,288 vertices and 15,483,485 unique directed edges) is reported as
a separate scalability endpoint. It is never included in a real-dataset
aggregate; the figures and cycle table place it after a dotted or ruled
separator.

\paragraph{SSSP interpretation.}
Spine starts from a verified persisted old-graph SSSP state and measures the
accepted update, automatic active-source discovery, incremental propagation,
and drain. The conversion-free GraSU+ReGraph baseline performs its declared
complete ReGraph execution because it has no equivalent persisted incremental
SSSP state. This is an architecture-level dynamic-service comparison, not an
identical-kernel microbenchmark.

\paragraph{Claim boundary.}
The frozen full matrix remains partial while expected rows are unfinished,
policy-stopped, or behavior-transition-invalidated; a row must complete or
receive a declared scope exclusion before the matrix can be called complete.
Device cycles, accepted memory bytes, and
DRAMSim3 HBM energy are simulator outputs. They do not claim cycle-for-cycle
FPGA calibration, on-chip dynamic energy, or total board power.

\paragraph{Behavior-transition coverage.}
The following 1 prior Spine rows are deliberately absent
until identical-case successors pass with identical final state: HW-SSSP.
They are not failures, zeros, or inputs to any aggregate.

\paragraph{Full-graph wall-time feasibility.}
An execution-feasibility model fitted to 3 completed,
correctness-admitted K4-shared SSSP rows projects total host times of SO 105.7 h, PK 342.2 h
at the observed cycle rates. Even a stress test with ten times less work and
twice the observed simulator rate leaves SO 3.6 h, PK 16.8 h, compared with the
3 h campaign budget. Those full-graph executions were therefore
soft-stopped by policy. Their partial cycles and memory requests are retained
only as host-runtime evidence and never enter accelerator-performance
aggregates. The validated source-matched preflights report LJ08 source 26, 35 supersteps, 1182.1 h, R19 source 113, 10 supersteps, 52.8 h, LJ source 71, 27 supersteps, 789.2 h. Their total-cycle estimates multiply validated oracle work by the median completed-run cycles/edge-round; they are projections, not completed cycle simulations. For the remaining 2 unlaunched real graphs, a conservative one-round screen uses the smallest completed-run cycles/edge-round and the largest completed-run simulator rate. Even the mandatory first full-graph round projects HW 35.7 h, OK 74.2 h; these screened rows were not launched and are not performance data.

\paragraph{Machine-checked evidence audit.}
All 25 observed executions pass
their individual architecture and mathematical correctness gates, closed AXI/HBM
arbitration and traffic ledgers, and DRAM request conservation. Component-level
activity is present for all observed executions
(163 normalized component rows).

\clearpage
\section{Current primary comparison}
\begin{figure}[H]
\centering
\includegraphics[width=0.98\linewidth]{\vfigdir/formal_v7_all_spine_e2e.pdf}
\caption{Absolute insertion-batch-8 device cycles for all
16 correctness-admitted current-version Spine rows. Blue
bars and unmarked orange bars are completed executions. Cross-hatched G+R bars
  marked Proj./T/O are total-cycle feasibility projections derived either from
  stopped execution prefixes or from validated source-matched host preflights;
  they are not measured performance. Upward triangles are conservative one-round
  feasibility screens, not total-cycle predictions. Downward triangles for
  SO CC and ResPR are strict monotonic device-cycle lower bounds from stopped
  incomplete executions; they are neither completed results nor projected totals.
The invalidated prior-version HW row and queued OK row have no admitted Spine
result and are therefore absent.}
\end{figure}

\noindent\textit{Projection method.}
For a source-matched preflight row, we compute
$C_{\mathrm{proj}}=\mathrm{median}_i
\{C_i/(E_iS_i)\}\,E_{\mathrm{target}}S_{\mathrm{target}}$.
The unit cost is fitted from three completed, correctness-admitted K4-shared
SSSP executions; the target edge count and minimum superstep count come from
the complete-graph host oracle using the same source cohort as Spine. These
values are feasibility projections, not completed simulations, and are
excluded from measured cross-architecture aggregates.

\begin{figure}[H]
\centering
\includegraphics[width=0.98\linewidth]{\vfigdir/formal_v7_primary_ratios.pdf}
\caption{Ratios for complete insertion-batch-8 pairs. Values above one favor
Spine for E2E latency and indicate that GraSU+ReGraph uses more memory traffic
or HBM energy in the other panels. R19, when complete, appears to the right of
the dotted separator and is not part of the real-dataset population.}
\end{figure}

\begin{table}[H]
\centering
\small
\input{\vdatadir/pair_table.tex}
\caption{Absolute device cycles. Only cross-system final-state-matched pairs
enter this table. R19 is separated from the real datasets by a rule.}
\end{table}

\clearpage
\section{Memory traffic and request locality}
\begin{figure}[H]
\centering
\includegraphics[width=0.98\linewidth]{\vfigdir/formal_v7_memory_locality.pdf}
\caption{Absolute accepted-backend traffic and request-stream locality for
the same complete pairs. A request is discontinuous when its accepted address
does not continue the previous request from the same initiator, operation, and
logical HBM channel. This is not a DRAM row-buffer-miss metric. R19 is a
separate synthetic endpoint.}
\end{figure}

\clearpage
\section{Structure-update throughput}
\begin{figure}[H]
\centering
\includegraphics[width=0.98\linewidth]{\vfigdir/formal_v7_update_throughput.pdf}
\caption{Correctness-gated AskUbuntu full-graph insertion throughput at
batch sizes $2^6$, $2^{10}$, $2^{14}$, and $2^{17}$. The current snapshot
contains 4 of 4 planned batch points. Spine counts maintenance through its
drained completion boundary; G+R counts GraSU PMA update completion. Both
exclude graph-algorithm iterations, equivalent to setting max-iter to zero.}
\end{figure}

\paragraph{Window boundary.}
This figure isolates graph-structure mutation. The counters come from dedicated
correctness-gated executions and stop at each architecture's serial
update-phase boundary. We use source vertex 0, which remains isolated after all
four batches, solely to shorten the post-update SSSP correctness drain. Both
architectures enforce an update--barrier--compute order, so this source choice
cannot affect the recorded update phase. For Spine the phase includes required
level selection, L0/carry work, metadata publication, and drain; it excludes
reader, propagation, convergence, and algorithm drain. For G+R it includes the
GraSU PMA update and excludes ReGraph execution. It must not be read as
end-to-end dynamic graph service latency.

\paragraph{Operation coverage.}
The separate insertion, deletion, and weight-change grid at batches 1, 8, and
64 remains in \texttt{update\_operation\_pairs.csv} and in the fallback/RQ3
evidence. It is not mixed into this scaling plot because non-monotonic SSSP
updates exercise a bounded full-rebuild policy rather than direct insertion.

\clearpage
\section{Simulator execution cost}
\begin{figure}[H]
\centering
\includegraphics[width=0.98\linewidth]{\vfigdir/formal_v7_simulator_runtime.pdf}
\caption{Host wall time and simulated-device-cycle throughput for the
correctness-matched pairs. These are simulator engineering diagnostics under
the recorded campaign concurrency, not accelerator latency or a
cross-architecture performance metric.}
\end{figure}

\paragraph{Use boundary.}
Wall time establishes experiment feasibility and identifies simulator
optimization targets. Architecture claims use device cycles and the shared
memory model; they never use host wall time.

\clearpage
\section{RQ3: realized work and latency}
\begin{figure}[H]
\centering
\includegraphics[width=0.98\linewidth]{\vfigdir/rq3_latency_breakdown.pdf}
\caption{Direct ten-stage, exclusive critical-path composition for zero-net,
shallow insertion, deep carry, residual PageRank correction, and SSSP/CC
nonmonotonic fallback. The six maintenance stages are counted in the execution
core; overlapping resolve/app intervals are assigned to the component gating
their completion. For residual PageRank, $T_{seed}$ also contains the serial
device correction/seed interval. Every bar closes exactly to E2E device cycles.}
\end{figure}

\begin{table}[H]
\centering
\small
\input{\rqdatadir/representative_table.tex}
\caption{Deterministically selected maximum-cycle representative of each
eligible realized-work class.}
\end{table}

\clearpage
\section{RQ3: cost-model validation}
\begin{figure}[H]
\centering
\includegraphics[width=0.98\linewidth]{\vfigdir/rq3_work_correlations.pdf}
\caption{Direct execution counters versus mechanism-local cycles across
95 dual-oracle-admitted executions. Blue squares are
calibration cases and red circles are holdout traces.}
\end{figure}

\begin{table}[H]
\centering
\small
\input{\rqdatadir/regression_table.tex}
\caption{Realized-work regressions. Carry work includes old payload reads,
merge inputs, and output rewrites; cursor inspection is reported separately.}
\end{table}

\paragraph{Interpretation.}
The seven panels report every requested realized-work relation, including weak
sort-only and drain-only relations rather than hiding them. The slope and $R^2$
table distinguishes mechanisms that transfer across mixed trace classes from
those that need a richer topology- or contention-aware predictor.

\begin{figure}[H]
\centering
\includegraphics[width=0.98\linewidth]{\vfigdir/rq3_e2e_cost_model.pdf}
\caption{Calibration-only nonnegative realized-work model versus measured E2E
cycles, plus absolute residuals. The holdout contains
39 real-trace executions that never participate in
fitting; holdout $R^2=0.933$, median absolute error
$26.5\%$, mean absolute error
$27.4\%$, and maximum absolute error
$72.0\%$.}
\end{figure}

\paragraph{RQ3 boundary.}
These correlations validate internal cost structure and bottleneck attribution;
they are not cycle-for-cycle FPGA calibration or proof that one scalar predicts
every topology. The timed device boundary starts with an already-resident sorted
update buffer: host DMA and an external FLiMS sort are not included. Current
$T_{seed}$ measures maintenance dirty-source publication and, for residual
PageRank, the serial device correction interval covering resident rank/degree
reads, physical-edge scans, seed writes, and active-list publication before
propagation. The host supplies only the immutable oracle/work trace.

\clearpage
\section{Implementation-level power attribution}
\begin{figure}[H]
\centering
\includegraphics[width=0.98\linewidth]{\vfigdir/formal_v7_component_power.pdf}
\caption{Vivado post-route vectorless hierarchy power at 150 MHz. The four
bars are distinct routed builds; the Spine evidence covers SSSP only. Values
use Vivado default activity with Low confidence and establish component
attribution, not workload-calibrated energy or board power.}
\end{figure}

\paragraph{Energy boundary.}
The workload-specific HBM energy ratios in Figure 1 come from DRAMSim3 and are
paired with the exact executions plotted there. The routed bars above establish
implementation-level component power but retain Vivado's Low-confidence
vectorless activity assumption.

\begin{figure}[H]
\centering
\includegraphics[width=0.98\linewidth]{\vfigdir/formal_v7_workload_energy.pdf}
\caption{Complete workload-duration energy ledger for the nine matched
SSSP, CC, and residual-PageRank pairs. DRAM energy comes from each exact
DRAMSim3 command stream. FPGA hierarchy components use execution-driven active
cycles multiplied by routed vectorless power; platform dynamic and device
static use the full measured interval. All ledgers close. This is a modeled
energy estimate, not RTL-SAIF or board-power measurement. Spine CC/Residual and
G+R CC use explicitly recorded routed-algorithm proxies.}
\end{figure}

\paragraph{Energy interpretation.}
This result has the same separation of concerns used throughout the report:
the cycle simulator supplies workload activity and duration, DRAMSim3 supplies
HBM-device energy, and implementation tools supply component power. CACTI
selected-SRAM ASIC projections remain a separate ledger and are not summed
with FPGA BRAM/URAM power. The largest remaining confidence upgrade is
algorithm-matched routed builds with RTL-derived SAIF activity.
\end{document}
